\subsection{对偶的函子性质}

设 G 与 H 是局部紧交换群。回顾：若 \(\varphi\colon\mathbf{G}\to\mathbf{H}\) 是拓扑群的一个态射，则对偶态射 \(\widehat{\varphi}\colon\widehat{\mathrm{H}}\to\widehat{\mathrm{G}}\) 由 \(\langle\chi,\varphi(x)\rangle=\langle\widehat{\varphi}(\chi),x\rangle\)（对一切 \(\chi\in\widehat{\mathrm{H}}\) 与 \(x\in\mathrm{G}\)）定义。这个定义表明 \(\widehat{\widehat{\varphi}}=\varphi\)，其中按照 II，第 220 页定理 2 把 G（相应地，H）与 \(\widehat{\widehat{\mathrm{G}}}\)（相应地，\(\widehat{\widehat{\mathrm{H}}}\)）等同。

设 \(\theta\) 是从 \(\mathbf{G} \times \mathbf{H}\) 到 \(\mathbf{U}\) 中的一个映射。对每个 \(x \in \mathbf{G}\)（相应地，每个 \(y \in \mathbf{H}\)），用 \(\theta_x\)（相应地，\(\theta^y\)）表示从 \(\mathbf{G}\) 到 \(\mathbf{U}\) 中的、由 \(y \mapsto \theta(x, y)\) 定义的函数（相应地，从 \(\mathbf{H}\) 到 \(\mathbf{U}\) 中的、由 \(x \mapsto \theta(x, y)\) 定义的函数）。假设映射 \(\alpha: x \mapsto \theta_x\) 是拓扑群 \(\mathbf{G}\) 到拓扑群 \(\widehat{\mathbf{H}}\) 上的一个同构。对每个 \(y \in \mathbf{H}\) 与 \(x \in \mathbf{G}\)，有

\[
\theta^ {y} (x) = \theta (x, y) = \langle \alpha (x), y \rangle = \langle x, \widehat {\alpha} (y) \rangle ,
\]

即 \(\theta^{y} = \widehat{\alpha}(y)\)。于是由 II，第 220 页定理 2，映射 \(\beta: y \mapsto \theta^{y}\) 是拓扑群 H 到拓扑群 \(\widehat{G}\) 上的一个同构。在这类条件下，我们说 \(\theta\) 使 G 与 H 成对偶，或说 G 与 H 关于 \(\theta\) 成对偶。这时我们将把 G 与 H 各自同另一个群的对偶等同起来。我们将把由 dx 的对偶测度经结构转移所得到的 H 上的 Haar 测度称为 Haar 测度 dx 的对偶测度。

引理 7.　设 \((\mathrm{G}_{i})_{i\in\mathrm{I}}\) 与 \((\mathrm{H}_{i})_{i\in\mathrm{I}}\) 是局部紧拓扑群的有限族。对 \(i\in I\)，设 \(\theta_{i}\colon G_{i}\times H_{i}\to U\) 是使群 \(G_{i}\) 与 \(H_{i}\) 成对偶的一个映射。映射 \(\theta\) 由

\[
\theta ((g _ {i}), (h _ {i})) = \prod_ {i \in I} \theta_ {i} (g _ {i}, h _ {i})
\]

定义，它使群 \(\prod G_{i}\) 与 \(\prod H_{i}\) 成对偶。

这由 II，第 206 页命题 2 及其前面的定义得到。

定义 5.　设 G、H 与 K 是拓扑群。设 \(f: \mathrm{H} \to \mathrm{G}\) 与 \(g: \mathrm{G} \to \mathrm{K}\) 是拓扑群的态射。若 \((f, g)\) 是群的一个正合列（A，II，第 10 页，注记 5），且 \(f\) 与 \(g\) 都是严格态射，则称它为拓扑群的一个正合列。

我们将用图式

\[
\mathrm{H} \xrightarrow {f} \mathrm{G} \xrightarrow {g} \mathrm{K},
\]

来表示一个正合列，并且我们说一个图式

\[
\mathrm{G} _ {1} \xrightarrow {f _ {1}} \mathrm{G} _ {2} \xrightarrow {f _ {2}} \mathrm{G} _ {3} \rightarrow \dots \rightarrow \mathrm{G} _ {n} \xrightarrow {f _ {n}} \mathrm{G} _ {n + 1}
\]

是正合的，如果每个二元组 \((f_i, f_{i+1})\)（\(1 \leqslant i \leqslant n-1\)）都是正合的。

一个序列

\[
1 \to \mathrm{H} \xrightarrow {f} \mathrm{G} \xrightarrow {g} \mathrm{K} \to 1
\]

是正合的，当且仅当 f 是一个严格单射态射，g 是一个严格满射态射，且 g 的核等于 f 的像。若 K 是分离的，则 f 的像是 G 的一个闭子群。

例.　1) 设 \(f\colon H \to G\) 是一个严格单射态射，其像是一个正规子群。序列

\[
1 \to \mathrm{H} \xrightarrow {f} \mathrm{G} \xrightarrow {p} \mathrm{G} / f (\mathrm{H}) \to 1,
\]

（其中 p 是典范投影）是正合的。特别地，若 H 是 G 的一个闭正规子群，则拓扑群的序列

\[
1 \to \mathrm{H} \xrightarrow {j} \mathrm{G} \xrightarrow {p} \mathrm{G} / \mathrm{H} \to 1,
\]

（其中 j 是包含映射，p 是典范投影）是正合的。

\begin{enumerate}
\def\labelenumi{\arabic{enumi})}
\setcounter{enumi}{1}
\tightlist
\item
  设 \(g: G \rightarrow K\) 是一个严格满射态射。序列
\end{enumerate}

\[
1 \to \mathrm{Ker} (g) \stackrel {j} {\longrightarrow} \mathrm{G} \stackrel {g} {\longrightarrow} \mathrm{K} \to 1,
\]

（其中 j 是包含映射）是正合的。

定理 4.　一个序列

\[
\mathrm{H} \xrightarrow {f} \mathrm{G} \xrightarrow {g} \mathrm{K}
\]

（为局部紧交换拓扑群的序列）是正合的，当且仅当对偶序列

\[
\widehat {\mathrm{K}} \xrightarrow {\widehat {g}} \widehat {\mathrm{G}} \xrightarrow {\widehat {f}} \widehat {\mathrm{H}}
\]

是正合的。

我们先证明几个引理。此外注意，它们中的每一个也都是定理 4 的断言的简单推论。

引理 8.　设 \(g \colon G \to K\) 是局部紧交换拓扑群的一个态射。若态射 \(g\) 是满射且严格的，则 \(\widehat{g}\) 是单射且严格的。

由于 \(g\) 是满射，态射 \(\widehat{g}\) 是单射（II，第 205 页，引理 2）。为证明 \(\widehat{g}\) 是一个严格态射，只需证明：对 \(e\) 在 \(\widehat{\mathbf{K}}\) 中的每个邻域 U，存在 \(e\) 在 \(\widehat{\mathbf{G}}\) 中的一个邻域 V，使得 \(\overline{\widehat{g}}^{-1} (\mathrm{V})\subset \mathrm{U}\)（II，第 200 页，引理 2）。设 U 是 \(e\) 在 \(\widehat{\mathbf{K}}\) 中的一个这样的邻域。由 \(\widehat{\mathbf{K}}\) 的拓扑的定义，存在 K 的一个紧子集 X 及一个数 \(\varepsilon >0\)，使得 U 包含由那些 \(\widehat{z}\in \widehat{\mathbf{K}}\)（它们对每个 \(z\in \mathrm{X}\) 都满足 \(|\langle \widehat{z},z\rangle -1| < \varepsilon\)）组成的集合。由于 \(g\) 严格且满射，由 TG，I，第 80 页，命题 10，存在 G 的一个紧子集 \(\mathrm{X}_0\)，使得 \(g(\mathrm{X}_0) = \mathrm{X}\)。设 V 是 \(e\) 在 \(\widehat{\mathbf{G}}\) 中的一个邻域，它由元素 \(\chi \in \widehat{\mathbf{G}}\)（它们对所有 \(x\in \mathrm{X}_0\) 都满足 \(|\langle \chi ,x\rangle -1| < \varepsilon\)）组成。于是有 \(\overline{\widehat{g}}^{-1} (\mathrm{V})\subset \mathrm{U}\)。这就证明了该断言。

引理 9.　设 \(f\colon H\to G\) 是局部紧拓扑群的一个态射。若态射 \(f\) 是单射且严格的，则 \(\widehat{f}\) 是满射且严格的。

假设 \(f\) 单射且严格。态射 \(\widehat{f}\) 通过过渡到商诱导一个态射 \(q \colon \widehat{\mathrm{G}} / \mathrm{Ker}(\widehat{f}) \to \widehat{\mathrm{H}}\)。剩下要证明这是一个拓扑群的同构；由对偶性，为此只需证明它的对偶 \(\widehat{q}\) 是一个同构。

用 L 表示 \(\widehat{\mathrm{G}}/\mathrm{Ker}(\widehat{f})\) 的对偶群，用 \(p\colon\widehat{\mathrm{G}}\to\widehat{\mathrm{G}}/\mathrm{Ker}(\widehat{f})\) 表示典范投影。我们先证明 \(\widehat{p}\) 通过过渡到子空间诱导一个从 L 到 \(f(\mathrm{H})\) 上的同构。

我们有 \(q \circ p = \widehat{f}\)，从而 \(\widehat{p} \circ \widehat{q} = f\)。因此 \(\widehat{p}\) 的像包含 \(f(\mathrm{H})\)。

由于 \(f\) 严格，它的像 \(f(\mathrm{H})\) 是 \(\mathrm{G}\) 的一个局部紧子群，因而是闭的（TG，III，第 22 页，推论 2）。设 \(\mathrm{K} = \mathrm{G}/f(\mathrm{H})\)，并考察正合序列

\[
1 \to \mathrm{H} \xrightarrow {f} \mathrm{G} \xrightarrow {g} \mathrm{K} \to 1
\]

相伴于例 1。由对偶性，态射 \(\widehat{f} \circ \widehat{g}\) 是平凡的，因而 \(\widehat{g}\) 的像含于 \(\operatorname{Ker}(\widehat{f}) = \operatorname{Ker}(p)\)。于是 \(p \circ \widehat{g}\) 是平凡态射，且再由对偶性，\(g \circ \widehat{p}\) 也是平凡的。由此推出 \(\widehat{p}\) 的像含于 \(g\) 的核，而后者等于 \(f(\mathrm{H})\)。我们断定 \(\widehat{p}\) 的像等于 \(f(\mathrm{H})\)。

此外，由于 p 是一个满射且严格的态射，对偶态射 \(\widehat{p}\) 是从 L 到 G 中的一个严格单射态射（引理 8）。由此，\(\widehat{p}\) 诱导一个从 L 到 \(f(\mathrm{H})\) 上的拓扑群同构。由于 \(\widehat{p} \circ \widehat{q} = f\)，且 f 诱导一个从 H 到其像 \(f(\mathrm{H})\) 上的同构，故态射 \(\widehat{q}\) 是一个同构。

引理 10.　设

\[
\mathrm{H} \xrightarrow {f} \mathrm{G} \xrightarrow {g} \mathrm{K}
\]

是局部紧交换群的一个正合序列。\(\hat{f}\) 的核等于 \(\hat{g}\) 的像。

同态 \(\widehat{f} \circ \widehat{g}\) 由对偶性是平凡的，故 \(\widehat{g}\) 的像含于 \(\widehat{f}\) 的核。反之，设 \(\chi\) 属于 \(\widehat{f}\) 的核。这意味着 \(\operatorname{Im}(f) = \operatorname{Ker}(g)\) 含于 \(\chi\) 的核，因而存在一个特征标 \(\eta\)（在 \(\operatorname{Im}(g)\) 上），使得 \(\eta \circ g = \chi\)。由于 \(\operatorname{Im}(g)\) 到 K 中的包含映射是严格的，从 K 到 \(\operatorname{Im}(g)\) 的特征标限制的对偶映射是满射的（引理 9）。因此存在 K 的一个特征标 \(\beta\)，使得 \(\eta\) 是 \(\beta\) 的限制，从而 \(\chi = \beta \circ g = \widehat{g}(\beta)\)。由此我们断定 \(\widehat{f}\) 的核含于 \(\widehat{g}\) 的像。

现在来证明定理 4。由对偶性，只需证明序列 \(\widehat{\mathrm{K}}\xrightarrow{\widehat{g}}\widehat{\mathrm{G}}\xrightarrow{\widehat{f}}\widehat{\mathrm{H}}\) 在序列 \(\mathrm{H}\xrightarrow{f}\mathrm{G}\xrightarrow{g}\mathrm{K}\) 正合时是正合的。而由引理 8 与引理 9，态射 \(\widehat{f}\) 与 \(\widehat{g}\) 都是严格的；由引理 10，\(\widehat{f}\) 的核等于 \(\widehat{g}\) 的像。

推论 1.　设

\[
1 \to \mathrm{H} \xrightarrow {f} \mathrm{G} \xrightarrow {g} \mathrm{K} \to 1
\]

是局部紧交换拓扑群的一个正合序列。态射 \(\widehat{g}\) 诱导 \(\widehat{\mathrm{K}}\) 与 \(f(\mathrm{H})^{\perp}\) 之间的一个同构，而 \(\widehat{f}\) 通过过渡到商诱导 \(\widehat{\mathrm{G}} / f(\mathrm{H})^{\perp}\) 与 \(\widehat{\mathrm{H}}\) 之间的一个同构。

由定理 4，序列

\[
1 \to \widehat {\mathrm{K}} \xrightarrow {\widehat {g}} \widehat {\mathrm{G}} \xrightarrow {\widehat {f}} \widehat {\mathrm{H}} \to 1\tag{32}
\]

是正合的。因此态射 \(\widehat{g}\) 诱导一个从 \(\widehat{\mathbf{K}}\) 到 \(\mathrm{Ker}(\widehat{f}) = f(\mathrm{H})^{\perp}\) 上的同构（II，第 205 页，引理 2），而 \(\widehat{f}\) 通过过渡到商诱导一个从 \(\widehat{\mathrm{G}} / \mathrm{Ker}(\widehat{f}) = \widehat{\mathrm{G}} / f(\mathrm{H})^{\perp}\) 到 \(\widehat{\mathrm{H}}\) 上的同构（同前）。

推论 2.　设 \(f\colon G\to H\) 是局部紧交换拓扑群的一个态射。态射 \(f\) 是严格的，当且仅当 \(\widehat{f}\) 是严格的。

由严格态射的典范分解（E，II，第 44 页），这可由引理 8 与引理 9 得到。

推论 3.　设 H 是 G 的一个子群。则 (H \(^{\perp}\) ) \(^{\perp}\) = \(\overline{H}\)。

由于 \(H^{\perp} = \overline{H}^{\perp}\)，我们可以假设 H 是闭的。设 \(f: H \to G\) 是典范单射，\(p: G \to G/H\) 是典范投影。我们有 \(H^{\perp} = \operatorname{Ker}(\widehat{f})\)（引理 10）。设 k 是 \(H^{\perp}\) 到 \(\widehat{G}\) 中的典范单射。由定理 4，态射 \(\widehat{p}\) 诱导一个拓扑群的同构 \(\widehat{p}_{H}: \widehat{G/H} \to H^{\perp}\)，且我们有 \(k \circ \widehat{p}_{H} = \widehat{p}\)。于是（推论 1）我们得到

\[
(\mathrm{H} ^ {\perp}) ^ {\perp} = \mathrm{Ker} (\widehat {k}) = \mathrm{Ker} (\widehat {\widehat {p}} _ {\mathrm{H}} \circ \widehat {k}) = \mathrm{Ker} (p) = \mathrm{H}.
\]

推论 4.　设 I 是一个集合，\((\mathrm{H}_{i})_{i\in\mathrm{I}}\) 是 G 的一族闭子群。由诸 \(H_{i}\) 生成的闭子群的正交是 \(\bigcap_{i\in I}H_{i}^{\perp}\)。\(\bigcap_{i}H_{i}\) 的正交是由诸子群 \(H_{i}^{\perp}\) 生成的闭子群。

第一个断言由正交的定义得到。把这个结果与推论 3 应用于闭子群族 \((\mathrm{H}_{i}^{\perp})_{i\in\mathrm{I}}\)（它们是 \(\widehat{G}\) 的闭子群），便知 \(\bigcap_{i}H_{i}\) 是由诸子群 \(H_{i}^{\perp}\) 生成的闭子群的正交，于是第二个断言由对偶性得到。

推论 5.　设 \(\varphi\) : G → H 是局部紧交换群的一个态射。则 H 的子群 \(\overline{\text{Im}(\varphi)}\) 与子群 \(\text{Ker}(\widehat{\varphi})\)（\(\widehat{\text{H}}\) 的子群）互为正交。特别地，\(\widehat{\varphi}\) 为单射的必要充分条件是 \(\varphi\) 的像在 H 中稠密。

我们有 \(\operatorname{Ker}(\widehat{\varphi}) = \varphi(\mathrm{G})^{\perp}\)（II，第 205 页，引理 2），由此据推论 3 即得结论。

推论 6.　设 \(k \in \mathbf{Z}\)。则从 G 到 G 的同态 \(x \mapsto x^k\) 的核与态射 \(\chi \mapsto \chi^k\)（从 \(\widehat{\mathrm{G}}\) 到 \(\widehat{\mathrm{G}}\)）的像的闭包互为正交。

由于从 G 到 G 的态射 \(x \mapsto x^{k}\) 与态射 \(\chi \mapsto \chi^{k}\)（从 \(\widehat{G}\) 到 \(\widehat{G}\)）互为对偶，这由上一个推论得到。

回顾（A，X，第 17 页）：若对每个非零的 \(n \in \mathbf{Z}\)，从 A 到 A 的映射 \(a \mapsto a^n\) 都是满射的，则称交换群 A 是可除的。

推论 7.　设 G 是一个局部紧交换群。

\begin{enumerate}
\def\labelenumi{\alph{enumi})}
\item
  若 G 是可除的，则对偶群 \(\widehat{\mathrm{G}}\) 是无挠的；
\item
  若对偶群 \(\widehat{\mathbf{G}}\) 是无挠的，且 \(k\in\mathbf{Z}\) 非零，则从 G 到 G 的同态 \(x\mapsto x^{k}\) 的像在 G 中稠密；
\item
  假设 G 是离散的或紧的。G 为可除的必要充分条件是 \(\widehat{G}\) 无挠。
\end{enumerate}

断言 \(a\)) 与 \(b\)) 由推论 6 得到。若 G 是离散的或紧的，则从 G 到 G 的同态 \(x \mapsto x^k\) 的像是闭的，而 \(c\)) 由 \(a\)) 与 \(b\)) 得到。

注.　存在不可除且使 \(\widehat{G}\) 无挠的局部紧交换群 G（II，第 299 页，习题 63）。

